% ═══════════════════════════════════════════════════════════════════════
\section{Sous quel angle regarder un nuage de points ?}
% ═══════════════════════════════════════════════════════════════════════

\begin{frame}{Un tableau de données est un nuage de points}
  \begin{itemize}
    \item Deux notes par élève (mathématiques, sciences) : chaque élève est
      un point du plan, la classe est un nuage de points.
    \item Trois mesures par manchot : un nuage de $\R^3$. Les dix épreuves
      d'un décathlon : $\R^{10}$. Des centaines de milliers de marqueurs
      génétiques : une dimension gigantesque.
    \item Dans le plan, un coup d'œil suffit pour repérer une tendance, des
      groupes, un point isolé.
    \item Mais nos yeux s'arrêtent à la dimension $3$ --- et sur une page ou
      un écran, nous ne voyons jamais qu'une projection sur un plan.
  \end{itemize}
\end{frame}

\begin{frame}{La photographie d'une théière}
  \begin{columns}[c]
    \begin{column}{0.55\textwidth}
      \begin{itemize}
        \item Photographiée de dessus : un disque.
        \item Photographiée de côté : on reconnaît le bec et l'anse.
        \item Deux projections du même objet ; l'une est beaucoup plus
          informative que l'autre.
      \end{itemize}
      \vspace{6pt}
      L'analyse en composantes principales (ACP) répond à la question
      correspondante pour un nuage de points :
      \begin{center}
        \color{sectioncolor}sous quel angle faut-il le regarder pour en
        voir le plus possible ?
      \end{center}
    \end{column}
    \begin{column}{0.42\textwidth}
      \FigBox[0.55\textheight]{%
        \begin{tikzpicture}
          % vue de dessus
          \draw[fig] (0,2.6) circle (0.8);
          \node[etiq, gray] at (0,1.55) {vue de dessus};
          % vue de côté
          \begin{scope}[yshift=-0.2cm]
            \draw[fig] (-0.8,0) .. controls (-0.9,0.9) and (0.9,0.9) .. (0.8,0)
              .. controls (0.6,-0.45) and (-0.6,-0.45) .. (-0.8,0) -- cycle;
            \draw[fig] (0.78,0.25) .. controls (1.2,0.4) and (1.35,0.7) .. (1.55,0.85)
              -- (1.5,0.95) .. controls (1.2,0.75) and (1,0.6) .. (0.72,0.5);
            \draw[steelblue4, thick] (-0.8,0.35) .. controls (-1.35,0.45) and (-1.35,-0.15) .. (-0.78,-0.12);
            \draw[fig] (-0.15,0.66) rectangle (0.15,0.8);
            \node[etiq, gray] at (0,-0.75) {vue de côté};
          \end{scope}
        \end{tikzpicture}}
    \end{column}
  \end{columns}
\end{frame}

\begin{frame}{Le chemin du chapitre}
  L'ACP cherche la droite, puis le plan, puis le sous-espace de petite
  dimension sur lequel la projection du nuage déforme le moins les données.
  La réponse : les vecteurs propres d'une matrice symétrique, la matrice
  de covariance.
  \vspace{6pt}
  \begin{enumerate}
    \item un peu de statistique descriptive (secondaire) ;
    \item six élèves : la droite qui représente le mieux un nuage du plan
      --- projection orthogonale et théorème de Pythagore ;
    \item les manchots de l'Antarctique : le meilleur plan dans $\R^3$ ;
    \item l'ACP en toute généralité, et le calcul avec Python ;
    \item un peu d'histoire ;
    \item au-delà de trois dimensions : le décathlon.
  \end{enumerate}
  \vspace{4pt}
  Deux sujets réunis : la statistique du secondaire et les vecteurs
  propres de l'annexe D.
\end{frame}
